%\baselineskip11.8pt plus.2pt minus.2pt


%\normalsize
%\vskip\baselineskip
\advance\baselineskip by-.2pt

\niebowspis{Niebo w~lutym}

\niebowtyt{Niebo w~lutym}

%\vspace*{-4pt}
{\begin{dwieszpalty}
\mathcode`\,="013B


Luty jest pierwszym z~czterech kolejnych miesięcy roku, gdy Słońce szybko wspina się na północ, przebywając każdego kolejnego tygodnia wyraźnie dłużej nad widnokręgiem. Do końca miesiąca zwiększy ono wysokość górowania o~ponad $8^{\circ}$, a~dzień wydłuży się prawie do~11~godzin.

Ekliptyka o~zmierzchu tworzy jeszcze większy kąt z~horyzontem niż w~styczniu, dzięki czemu na początku lutego bardzo dobrze widoczny jest Księżyc tuż po nowiu oraz znajdujące się na wschód od Słońca i~niezbyt daleko od niego planety Układu Słonecznego. Srebrny Glob zacznie miesiąc w~fazie $12\%$ między planetami Saturn i~Wenus, około $4^{\circ}$ od drugiej z~nich. Saturn zbliża się do marcowej koniunkcji ze Słońcem, dlatego można go dostrzec tylko w~pierwszej części miesiąca. Potem zaginie w~zorzy wieczornej. A~szkoda, gdyż równonoc na Saturnie i~przejście Ziemi przez płaszczyznę jego pierścieni są coraz bliżej. 8~lutego od zmierzchu do zachodu planety można jeszcze spróbować zaobserwować przejście cienia Tytana po jej tarczy. Saturn świeci z~jasnością $+1,1^m$, przy średnicy tarczy~$16''$.

Wenus również zmniejsza odległość kątową do Słońca po styczniowej maksymalnej elongacji, pozostanie jednak ozdobą wieczornego nieba do końca zimy. W~lutym druga planeta od Słońca jest widoczna bardzo dobrze, zajmując o~zmierzchu pozycję na wysokości przekraczającej~$20^{\circ}$. Początkowo Wenus wędruje wzdłuż ekliptyki, ale pod koniec lutego zacznie wyraźnie odbijać na północ, szykując się do marcowej koniunkcji dolnej ze Słońcem. Planeta szybko zbliża się do nas, w~związku z~czym jej tarcza rośnie, zmniejszając jednocześnie swoją fazę. Na początku miesiąca jej tarcza ma średnicę $32''$ i~fazę $38\%$, by do końca lutego jej średnica zwiększyła się do $49''$, faza zaś zmniejszy się do $14\%$. Świeci przy tym bardzo jasno, przekraczając blask $-4,5^m$. To oznacza, że jest ona atrakcyjnym celem nawet dla właścicieli lornetek.
%\vspace*{\parskip}


Srebrny Glob dotrzyma towarzystwa Wenus jeszcze 2~lutego, pokazując się w~fazie $21\%$ jakieś $10^{\circ}$ na wschód od niej, a~następnie powędruje w~kierunku Plejad, które minie w~nocy z~5 na 6~lutego. Do tego czasu jego faza zwiększy się do I~kwadry. Tym razem Europa ma pecha, gdyż Księżyc zakryje niektóre gwiazdy gromady Plejad, będąc już pod horyzontem. W~Polsce 5~lutego Księżyc zdąży zbliżyć się do Plejad na $5^{\circ}$, a~następnej doby o~zmierzchu pokaże się również $5^{\circ}$, ale na wschód od nich. Tej nocy naturalny satelita Ziemi zbliży się także do Jowisza. Tuż przed zachodem obu ciał niebieskich dystans między nimi zmniejszy się także do $5^{\circ}$. Największa planeta Układu Słonecznego w~lutym jest widoczna wciąż całkiem dobrze, przecinając południk lokalny na początku miesiąca na całkiem ciemnym niebie. Jowisz wędruje niecałe $5^{\circ}$ na północ od Aldebarana, najjaśniejszej gwiazdy Byka, a~w~trakcie miesiąca jego jasność spadnie do $-2,3^m$, średnica tarczy natomiast do $40''$.

W~nocy z~7 na 8~lutego Księżyc minie gwiazdę El Nath, drugą co do jasności gwiazdę Byka, zwiększając przy tym fazę do $78\%$. U~nas minimalna odległość między tymi ciałami niebieskimi zmniejszy się do~około~$45'$. Dwie noce później zaś Srebrny Glob przeniesie się do wschodniej części Bliźniąt, gdzie odwiedzi najpierw planetę Mars, a~następnie Kastora i~Polluksa. Wieczorem Księżyc zbliży się do Marsa na mniej niż $0,5^{\circ}$, a~nad ranem zmniejszy dystans do Polluksa do $3^{\circ}$. Sam Mars szybko oddala się od Ziemi i~jego tarcza wyraźnie traci na średnicy i~jasności. Planeta prawie do końca miesiąca porusza się ruchem wstecznym, oddalając się od Polluksa na ponad~$7^{\circ}$. W~lutym jasność Marsa spadnie z~$-1^m$ do~$-0,3^m$, a~średnica jego tarczy -- z~$14''$ do $11''$. Podczas górowania Mars wznosi się na ponad $65^{\circ}$.

12 lutego nastąpi pełnia Księżyca i~jednocześnie zbliży się on na $1,5^{\circ}$ do Regulusa, najjaśniejszej gwiazdy Lwa. 5~nocy później faza Srebrnego Globu spadnie do $82\%$ i~dotrze on do środkowej części Panny, wędrując jakieś $5^{\circ}$ od Spiki, najjaśniejszej gwiazdy tej konstelacji. 20 lutego Księżyc osiągnie ostatnią kwadrę, a~dobę później wzejdzie w~gwiazdozbiorze Skorpiona, niecałe $3^{\circ}$ od Antaresa, jego najjaśniejszej gwiazdy. Potem Księżyc podąży ku nowiu, przez który przejdzie 28 lutego, ale ze względu na niekorzystne nachylenie ekliptyki i~dodatkowo przebywanie na południe od niej zniknie on w~zorzy porannej już kilka dni wcześniej.

Przez cały miesiąc bardzo dobrze widoczna jest gwiazda R~Leonis, której maksimum aktywności przewiduje się na 4~lutego. Jest to miryda, czyli długookresowa gwiazda zmienna, której jasność waha się od $+4,4^m$ do $+11^m$ w~okresie 312 dni. Za każdym razem zarówno w~maksimum, jak i~w~minimum blasku jej jasność może być różna, ale jeśli zbliży się ona do swojej największej obserwowanej jasności, to na ciemnym niebie da się ją łatwo dostrzec gołym okiem. R~Leo znajduje się jakieś $5^{\circ}$ na zachód od Regulusa, niedaleko gwiazd 6. wielkości 18~i~19 Leonis, z~którymi można porównywać jej blask. Ma też wyraźnie widoczną wiśniową barwę. Łatwo ją zatem zidentyfikować przez lornetkę. W~lutym gwiazda góruje przed północą na wysokości $50^{\circ}$.

\medskip
\rightline{\large
  \it  Ariel MAJCHER}
\end{dwieszpalty}

}\eject
\normalsize